
\فصل{تصویر موفقیت}

\قسمت{تعریف موفقیت}

ساکنا زمانی موفق است که:

\شروع{شمارش}
\فقره تمام کاربران (ساکنین، مدیر، کارکنان) بتوانند بدون آموزش تخصصی از سامانه استفاده کنند.
\فقره فرآیند پرداخت شارژ که قبلاً چند روز طول می‌کشید، در کمتر از ۵ دقیقه انجام شود.
\فقره ۱۰۰ درصد درخواست‌های تعمیراتی در سامانه ردیابی‌پذیر باشند و هیچ درخواستی گم نشود.
\فقره رزرو امکانات مشترک بدون تداخل و اشتباه انجام شود.
\فقره مدیر ساختمان زمان کمتری را صرف امور اداری تکراری کند.
\فقره مناقشات مالی بین مدیر و ساکنین به حداقل برسد.
\پایان{شمارش}

\قسمت{شاخص‌های کلیدی موفقیت}

\begin{table}[ht]
\centering
\caption{شاخص‌های کلیدی موفقیت (\lr{KPI})}
\label{جدول:KPI}
\begin{tabular}{|>{\raggedleft\arraybackslash}p{5cm}|>{\raggedleft\arraybackslash}p{2.5cm}|>{\raggedleft\arraybackslash}p{4.5cm}|}
\hline
\rl{\textbf{شاخص}} & \rl{\textbf{هدف}} & \rl{\textbf{روش اندازه‌گیری}} \\
\hline
\rl{نسبت پرداخت آنلاین شارژ} & \rl{بیش از ۸۰\%} & \rl{تعداد پرداخت‌های آنلاین / کل شارژها} \\
\hline
\rl{درخواست‌های خدماتی ردیابی‌شده} & \rl{۱۰۰\%} & \rl{تعداد درخواست‌های ثبت‌شده در سامانه} \\
\hline
\rl{زمان پاسخ سامانه} & \rl{کمتر از ۳ ثانیه} & \rl{میانگین زمان پاسخ \lr{API}} \\
\hline
\rl{رضایت مدیر از کاهش کار دستی} & \rl{مثبت} & \rl{بازخورد کیفی مدیر در دمو} \\
\hline
\rl{عدم تداخل رزرو} & \rl{صفر تداخل} & \rl{تعداد رزروهای متداخل ثبت‌شده} \\
\hline
\end{tabular}
\end{table}

\قسمت{تصویر موفقیت در پایان پروژه}

در پایان ایتریشن دوم (تیر ۱۲)، انتظار داریم:

\شروع{فقرات}
\فقره یک نسخه‌ی کاربردی از ساکنا آماده و قابل دمو باشد که تمام قابلیت‌های \lr{Must Have} پیاده‌سازی شده باشند.
\فقره سیستم احراز هویت، مدیریت واحدها، درخواست‌های خدماتی، پرداخت شارژ، رزرو امکانات و اطلاعیه کار کنند.
\فقره استاد درس بتواند کاربرد صحیح متدولوژی \lr{Lean DAD} و مستندات آن را در پروژه مشاهده کند.
\فقره تیم مهارت‌های لازم برای چرخه‌ی \lr{CD Lean} را کسب کرده باشد.
\پایان{فقرات}
